% !TeX root = ../main.tex
\section{Queue-Depth Sweep}\label{sec:queue-depth}

\begin{hypothesis}
  IOPS should rise steeply with queue depth as the SSD's parallel NAND channels fill, then saturate at a
  knee beyond which latency grows roughly linearly with queue depth for no extra throughput. By Little's
  Law the number of outstanding I/Os $L=\lambda W$ should track the configured iodepth while the queue
  stays full; where it stops tracking reveals a host-side or engine limit.
\end{hypothesis}

\begin{knobs}
  Fixed       & \qty{4}{\kibi\byte} random read; one job; \code{O_DIRECT} \\
  Varied      & iodepth \numrange{1}{128} \\
  Effective QD & verified in-band via the Little's Law table below (no \code{iostat} available) \\
  Script      & \code{SSD_Performance_Profiling/scripts/run.sh queue_depth} \\
\end{knobs}

\begin{figure}[tbp]
  \centering
  \IfDataExists{queue_depth}{%
    \begin{tikzpicture}
      \begin{axis}[acs, width=0.78\linewidth, xlabel={IOPS}, ylabel={Latency (\si{\micro\second})},
                   legend pos=north west, nodes near coords, point meta=explicit symbolic,
                   every node near coord/.append style={font=\scriptsize, anchor=south east}]
        \addplot+ table[x=iops_median, y=lat_mean_us_median, meta=qd] {data/queue_depth.csv};
        \addlegendentry{mean latency}
        \addplot+ table[x=iops_median, y=lat_p99_us_median, meta=qd] {data/queue_depth.csv};
        \addlegendentry{p99 latency}
      \end{axis}
    \end{tikzpicture}%
  }
  \caption{Throughput vs.\ latency; each point is labelled with its queue depth. IOPS saturates near
  \qty{158}{k} at QD=64; past the knee, QD=128 doubles latency for no extra throughput.}\label{fig:queue-depth}
\end{figure}

\subsection{Little's Law check}
With the queue kept full the outstanding-I/O count should satisfy $L = \lambda W$ (IOPS $\times$ mean
latency) and equal the configured iodepth \cite{little1961proof}.

\begin{table}[tbp]
  \centering\small
  \caption{Little's Law: measured $\lambda W$ against the configured queue depth
  (\qty{4}{\kibi\byte} random read).}\label{tab:littles-law}
  \begin{tabularx}{\linewidth}{@{}rrrrX@{}}
    \toprule
    iodepth & $\lambda$ (IOPS) & $W$ (\si{\micro\second}) & $L=\lambda W$ & \\
    \midrule
    1   & \num{20137}  & \num{49.6}  & \num{1.00}   & \\
    2   & \num{26325}  & \num{75.9}  & \num{2.00}   & \\
    4   & \num{30416}  & \num{131.3} & \num{3.99}   & \\
    8   & \num{46398}  & \num{172.1} & \num{7.99}   & \\
    16  & \num{74420}  & \num{213.9} & \num{15.92}  & \\
    32  & \num{125097} & \num{253.7} & \num{31.74}  & knee region \\
    64  & \num{158064} & \num{403.2} & \num{63.73}  & saturation \\
    128 & \num{158551} & \num{804.8} & \num{127.6}  & latency-only growth \\
    \bottomrule
  \end{tabularx}
\end{table}

\subsection{Analysis}
IOPS climbs almost linearly through QD=16 (\num{20}k\,$\to$\,\num{74}k), begins to bend at QD=32
(\num{125}k) and saturates at QD=64 (\num{158}k); QD=128 adds \num{0.3}\% throughput while mean latency
doubles from \qty{403}{} to \qty{805}{\micro\second}. The knee is therefore at QD\,$\approx$\,64, where the
device's channel parallelism is fully engaged. The Little's Law column is strikingly clean: $\lambda W$
equals the configured iodepth to within a percent at \emph{every} depth including QD=128
($L=\num{127.6}$), so the \code{libaio} engine really does keep the queue full---the effective queue depth
matches the requested one, which is the in-band substitute for an \code{iostat} aqu-sz check. This is also
the textbook signature of saturation: past the knee $\lambda$ is flat, so the only way $L$ can keep
growing with iodepth is for $W$ to grow proportionally---exactly the linear latency rise we see. The
practical reading: QD\,$\approx$\,32--64 extracts essentially all the throughput at a quarter-millisecond
latency; deeper queues only buy latency. Here the saturation is the device's, not the harness's, since
Little's Law holds all the way out.
